\chapter{Background}
\label{ch:background}

\section{Document Image Analysis}

Document image analysis (\ac{DIA}) is the field concerned with automatically extracting
structured information from images of documents~\cite{nagy2000twenty}. Unlike natural
scene analysis, document images have a strong spatial grammar: elements are arranged
according to a layout that carries semantic meaning. The position of a field, its font weight,
its proximity to a barcode—all of these encode information that a purely text-based approach
would lose.

For IKEA labels specifically, the document is constrained and controlled: the label designer
works from a template that dictates exactly which zone each element occupies. This makes
labels a particularly tractable sub-problem within \ac{DIA}: the vocabulary of layouts is
finite and enumerable, unlike free-form documents such as contracts or letters.

\subsection{Zone-Based Layout Analysis}

A canonical approach in \ac{DIA} is to segment the document into zones—rectangular or
polygonal regions—before performing per-zone recognition~\cite{breuel2003pages}. IKEA
labels are divided into named zones such as \texttt{left\_identity}, \texttt{center\_barcode},
\texttt{right\_compliance}, and \texttt{right\_date}, each with a defined stacking direction
and set of elements. Zone segmentation greatly reduces the search space for field extraction
and can be anchored to detected barcodes, whose position is both reliably detectable and
semantically significant.

\subsection{PDF-Native vs.\ Raster Extraction}

When a label is available as a PDF generated programmatically (e.g.\ from a design tool
such as Adobe Illustrator or a label template system), the text spans, font metadata, and
element positions are embedded directly in the PDF vector stream. PyMuPDF (fitz) can
extract these spans with sub-millimetre precision and confidence~1.0, entirely bypassing the
need for \ac{OCR}. This is denoted \emph{Pass~1} in the pipeline.

When only a raster image is available—a photograph, scan, or JPG export—\ac{OCR} becomes
necessary. The pipeline then runs \emph{Pass~5} (OCR ensemble), which is substantially
less reliable than PDF-native extraction, motivating the effort devoted to OCR engine
selection and ensemble design in \cref{ch:implementation}.

\section{Optical Character Recognition}

\ac{OCR} is the task of converting a raster image of text into a machine-readable character
sequence. The field has evolved from template-matching approaches through Hidden Markov
Models to the current generation of deep learning engines that jointly model image features
and language context~\cite{shi2016end}.

\subsection{Tesseract}

Tesseract is an open-source \ac{OCR} engine originally developed by HP and now maintained
by Google~\cite{smith2007tesseract}. Version~5 replaced the earlier LSTM architecture with
a fully convolutional LSTM network trained on a large multilingual corpus. Tesseract performs
particularly well on clean, high-resolution images with consistent typography. It is invoked
in the pipeline with the \texttt{pytesseract} Python binding.

\subsection{EasyOCR}

EasyOCR~\cite{easyocr2020} is a Python library implementing a CRAFT-based text detector
followed by a CRNN sequence recogniser. It handles curved and rotated text better than
Tesseract and is especially strong on non-Latin scripts. EasyOCR produces per-word
bounding boxes and confidence scores, which the pipeline uses for confidence gating
(words below~0.25 confidence are excluded from field mapping).

\subsection{PaddleOCR}

PaddleOCR (PP-OCRv5) from Baidu~\cite{du2020pp} is the state-of-the-art open-source
\ac{OCR} system as of 2024--2025 and was evaluated during the project. Platform
incompatibilities between PaddlePaddle and macOS arm64 / Python~3.13 prevented its
inclusion in the primary evaluation; however, benchmarking results on a Linux auxiliary
machine showed PaddleOCR achieving character-level accuracy of~92.3\% on the label
dataset, suggesting it would be the preferred engine for a Linux production deployment.

\subsection{OCR Ensemble Strategies}

Combining multiple \ac{OCR} engines is a well-established technique for improving
robustness~\cite{ratzlaff2003method}. The most common strategies are voting (majority
decision on character level), confidence weighting (weighted average of engine-assigned
confidences), and late fusion (merge detected bounding boxes by IoU, then pick the
higher-confidence word). This project implements late-fusion merging, where overlapping
word boxes from Tesseract and EasyOCR are merged using intersection-over-union (\ac{SSIM})
thresholding and the higher-confidence word is retained.

\section{Barcode Symbologies}

IKEA labels employ three barcode types, each encoding different layers of product information.

\subsection{ITF-14}

Interleaved Two-of-Five (ITF-14) is a linear barcode that encodes a 14-digit number—typically
a GTIN-14 logistics identifier—by interleaving two characters per symbol pair. IKEA uses
ITF-14 at various magnification levels (25\%--100\%) and with a mandatory 10-module quiet
zone on each side. The pyzbar library decodes ITF-14 by calling the underlying ZBar C library,
which operates on a grayscale pixel projection of the barcode area.

\subsection{DataMatrix}

DataMatrix is a 2D matrix barcode capable of encoding up to 3116 numeric or 2335 alphanumeric
characters in a compact footprint. IKEA uses DataMatrix exclusively in GS1 format,
encoding application identifiers (AIs) such as AI(240) for the item product code, AI(13)
for the packaging date, and AI(10) for the batch/lot number~\cite{gs12023general}.
DataMatrix is decoded using the pylibdmtx library, a Python binding of libdmtx.

\subsection{EAN-13}

EAN-13 (European Article Number) is the standard consumer product barcode, encoding a
13-digit GTIN-13. IKEA uses EAN-13 at 80\% and 100\% magnification. EAN-13 is decoded
via pyzbar alongside ITF-14.

\section{Rule-Based Validation}

Rule-based systems encode domain knowledge as explicit logical conditions that can be
evaluated deterministically. For document validation, rules typically cover:

\begin{itemize}
  \item \textbf{Field presence}: is a required element present in the extraction?
  \item \textbf{Format rules}: does the value conform to a regular expression or enumerated
  set?
  \item \textbf{Cross-field consistency}: do two fields that encode related information agree?
  \item \textbf{Structural rules}: are elements in the expected spatial zones?
\end{itemize}

Rule-based systems are inherently interpretable: every failure can be attributed to a specific
rule, with the extracted value and the expected value presented side-by-side. This is a key
advantage over neural classification approaches in a compliance context, where human
reviewers must be able to understand and act on the system's output.

\subsection{JSON Schema-Driven Rules}

In this project, validation rules are encoded as JSON objects within schema files—one file
per label type. Each rule specifies a \texttt{field}, a \texttt{check} type (\texttt{present},
\texttt{regex\_match}, \texttt{exact\_match}, \texttt{ai\_field\_present}), a
\texttt{severity} (\texttt{error} or \texttt{warning}), and the expected value or pattern.
This schema-driven design makes it straightforward to add new label types by writing a JSON
file, without modifying any Python code.

\section{Explainability in Automated Validation}

\ac{XAI} is the study of methods that make the outputs of \ac{ML} systems understandable
to humans~\cite{adadi2018peeking}. In the context of label validation, explainability is
not merely a desirable property—it is a functional requirement. A QA operator receiving a
``non-compliant'' decision must know:

\begin{enumerate}
  \item \textit{Which} field triggered the failure.
  \item \textit{What} value was found versus what was expected.
  \item \textit{Where} on the label the problematic element is located.
\end{enumerate}

The rule-based nature of the validator in this work provides points~1 and~2 by construction:
each \texttt{RuleResult} object carries the rule ID, the extracted value, the expected value,
and a human-readable message. Point~3 is addressed by attaching bounding-box coordinates
from the \ac{OCR} engine to each extracted field, enabling the \ac{UI} to highlight the
relevant region on the label image.

\section{GS1 Application Identifier Parsing}

GS1 Application Identifiers are two-to-four-digit numeric prefixes that identify the semantics
of the data that follows them in a GS1-encoded barcode~\cite{gs12023general}. For example:
\begin{itemize}
  \item \textbf{AI(240)}: Item product code (up to 30 characters).
  \item \textbf{AI(11)}: Production date, in YYMMDD format.
  \item \textbf{AI(13)}: Packaging date, in YYMMDD format.
  \item \textbf{AI(10)}: Batch or lot number.
\end{itemize}

Parsing GS1 content requires handling both the legacy bracket-encoded form
(\texttt{(240)12345678\ldots}) and the compact concatenated form
(\texttt{24012345678\ldots}) used in printed DataMatrix symbols. The pipeline implements
a GS1 parser that recognises both forms, handles variable-length AIs terminated by
function code~1 (FNC1) delimiters, and validates date fields against format constraints.
